\documentclass[aspectratio=169,professionalfonts]{ctexbeamer}

\usetheme{Boadilla}
\useinnertheme{rectangles}
\usefonttheme{professionalfonts}
\setbeamertemplate{navigation symbols}{}

\usepackage{amsmath}
\usepackage{amssymb}
\usepackage{bm}
\usepackage{booktabs}
\usepackage{mathtools}

\definecolor{DeepBlue}{RGB}{20,45,110}
\definecolor{SoftBlue}{RGB}{232,238,250}
\definecolor{Accent}{RGB}{180,60,40}
\definecolor{DarkText}{RGB}{35,35,35}

\setbeamercolor{title}{fg=white,bg=DeepBlue}
\setbeamercolor{frametitle}{fg=DeepBlue,bg=SoftBlue}
\setbeamercolor{structure}{fg=DeepBlue}
\setbeamercolor{normal text}{fg=DarkText,bg=white}
\setbeamercolor{block title}{fg=white,bg=DeepBlue}
\setbeamercolor{block body}{fg=DarkText,bg=SoftBlue}
\setbeamercolor{item projected}{fg=white,bg=DeepBlue}

\setbeamertemplate{title page}{
  \vspace{1.2cm}
  \begin{beamercolorbox}[wd=\paperwidth,ht=2.3cm,dp=0.4cm,leftskip=1cm]{title}
    {\LARGE \inserttitle\par}
    \vspace{0.35cm}
    {\large \insertauthor\par}
    \vspace{0.15cm}
    {\normalsize \insertdate\par}
  \end{beamercolorbox}
}

\newcommand{\R}{\mathbb{R}}
\newcommand{\Loss}{\mathcal{L}}
\newcommand{\Stop}{\operatorname{sg}}
\newcommand{\MSE}{\operatorname{MSE}}
\newcommand{\softmaxop}{\operatorname{SoftSort}}
\newcommand{\Sinkhorn}{\operatorname{Sinkhorn}}
\newcommand{\argsort}{\operatorname{argsort}}

\title{基于 OT 松弛的 Channel Grouping 优化}
\author{组会汇报}
\date{\today}

\begin{document}

\begin{frame}
  \titlepage
\end{frame}

\begin{frame}{1. 问题背景}
  考虑 MoE 中某个 expert 的 \texttt{down\_proj}：
  \[
    x \in \R^{B \times D},\qquad
    w \in \R^{O \times D},\qquad
    y_{\mathrm{ref}} = x w^\top \in \R^{B \times O}
  \]

  \vspace{0.4em}
  在不做通道重排时，直接量化得到
  \[
    y_q^{(\mathrm{id})} = Q(x)\,Q(w)^\top
  \]
  其重建误差通常受输入通道排列方式影响较大，特别是在分组量化场景下：
  \[
    \Loss_{\mathrm{id}} = \MSE\!\left(y_q^{(\mathrm{id})},\, y_{\mathrm{ref}}\right)
  \]

  \vspace{0.4em}
  本工作的核心问题是：
  \begin{block}{目标}
    在不修改权重数值的前提下，只通过学习 channel grouping，降低 MXFP4 分组量化误差。
  \end{block}
\end{frame}

\begin{frame}{2. 核心思路}
  真实关心的对象不是完整 permutation，而是 grouping：
  \[
    \mathcal{G} = \{G_1,\dots,G_M\}, \qquad |G_m| = k,\qquad M=D/k
  \]

  \vspace{0.4em}
  但为了保持可导训练，我们不直接优化 grouping，而是：
  \begin{enumerate}
    \item 用一维 score 诱导排序
    \item 用 Sinkhorn 得到位置级 soft assignment $P_{\mathrm{soft}}$
    \item 再用固定块矩阵把位置聚合成组
  \end{enumerate}

  \vspace{0.4em}
  因而我们保留了 OT 风格的连续松弛，同时把最终语义切回 group。
\end{frame}

\begin{frame}{3. 参数化：从离散 permutation 到连续 score}
  直接优化离散 permutation 很困难，因此引入一维可学习参数
  \[
    s = (s_1,\dots,s_D) \in \R^D
  \]

  \vspace{0.4em}
  先做标准化：
  \[
    \tilde{s} = \frac{s-\operatorname{mean}(s)}{\operatorname{std}(s)+\varepsilon}
  \]

  \vspace{0.4em}
  再由 score 诱导排序。离散路径上，
  \[
    \pi_{\mathrm{hard}} = \argsort(\tilde{s})
  \]
  并进一步构造 hard permutation matrix
  \[
    P_{\mathrm{hard}}[r,\pi_{\mathrm{hard}}(r)] = 1
  \]

  \vspace{0.4em}
  这样，每个 expert 的训练参数仅为长度 $D$ 的 score 向量，参数量非常小。
\end{frame}

\begin{frame}{4. 位置级 OT 松弛}
  对标准化后的 score，先排序得到
  \[
    \tilde{s}_{(1)} \le \tilde{s}_{(2)} \le \cdots \le \tilde{s}_{(D)}
  \]

  构造 cost matrix：
  \[
    C_{r,i} =
    \begin{cases}
      |\tilde{s}_{(r)}-\tilde{s}_i|, & \text{L1 metric}\\[0.3em]
      (\tilde{s}_{(r)}-\tilde{s}_i)^2, & \text{L2 metric}
    \end{cases}
  \]

  再定义
  \[
    \log \alpha = -\frac{C}{\tau}
  \]
  通过 Sinkhorn 归一化得到 soft position assignment：
  \[
    P_{\mathrm{soft}} = \Sinkhorn(\exp(\log \alpha))
  \]

  其可看作熵正则 OT / assignment：
  \[
    \min_{P \in \mathcal{B}}
    \langle P, C \rangle - \tau H(P)
  \]
  其中 $\mathcal{B}$ 是双随机矩阵集合。

  \vspace{0.3em}
  其中温度 $\tau$ 控制软硬程度：
  \[
    \tau \downarrow \quad \Longrightarrow \quad P_{\mathrm{soft}} \text{ 更接近一一匹配}
  \]
\end{frame}

\begin{frame}{5. 从位置到分组：\texorpdfstring{$A_{\mathrm{soft}} = P_{\mathrm{soft}}^\top B$}{Asoft = Psoft^T B}}
  定义固定块矩阵
  \[
    B \in \{0,1\}^{D \times M}, \qquad
    B[r,m]=1 \Leftrightarrow r \text{ 属于第 } m \text{ 个组的位置块}
  \]

  于是 group assignment 写成：
  \[
    A_{\mathrm{soft}} = P_{\mathrm{soft}}^\top B \in \R^{D \times M}
  \]

  展开后：
  \[
    A_{\mathrm{soft}}[i,m]
    =
    \sum_{r=1}^{D} P_{\mathrm{soft}}[r,i]\, B[r,m]
  \]

  含义：
  \[
    \text{channel } i \text{ 落入第 } m \text{ 个 group 的软权重}
  \]

  这一步等价于把 position-level OT plan 做 block aggregation。
\end{frame}

\begin{frame}{6. Group-level STE}
  由 hard 排序同样可得到 hard grouping：
  \[
    A_{\mathrm{hard}} = P_{\mathrm{hard}}^\top B
  \]

  训练中使用 group-level STE：
  \[
    A_{\mathrm{ste}}
    = A_{\mathrm{hard}} - \Stop(A_{\mathrm{soft}}) + A_{\mathrm{soft}}
  \]

  为了在代码中保持前向可执行，我们进一步把 soft grouping 均匀摊回位置空间：
  \[
    P_{\mathrm{group\_soft}} = \frac{1}{k} B A_{\mathrm{soft}}^\top
  \]

  \begin{block}{解释}
    前向仍使用 hard 排序对应的真实分组；
    反向梯度则通过 group-level soft surrogate 传播，从而让训练对象真正变成 grouping。
  \end{block}
\end{frame}

\begin{frame}{7. MXFP4 Fake Quant 与量化 STE}
  对重排后的张量按最后一维分组，group size 记为 $g$。对每个分组：
  \[
    a_{\max} = \max |x|
  \]
  主 scale 取为
  \[
    s_{\mathrm{blk}} = 2^{\left\lfloor \log_2\left(\frac{8}{7} a_{\max}\right) \right\rfloor}
  \]

  再对归一化值计算 secondary scale，并执行离散化。最终得到 fake quant 值 $Q(x)$。

  \vspace{0.4em}
  为了让量化可训练，使用
  \[
    Q_h(x) = x + \Stop(Q(x)-x)
  \]

  因而：
  \[
    \text{forward 用 } Q(x),\qquad
    \text{backward 近似看作 identity}
  \]

  当前方案中存在两层连续近似：
  \[
    s \rightarrow P_{\mathrm{soft}} \rightarrow A_{\mathrm{soft}} \rightarrow Q_h(\cdot)
  \]
\end{frame}

\begin{frame}{8. 损失函数}
  参考输出定义为
  \[
    y_{\mathrm{ref}} = xw^\top
  \]

  训练时实际执行的前向路径仍可写成
  \[
    x_{\mathrm{perm}} = xP_{\mathrm{hard}}^\top,\qquad
    w_{\mathrm{perm}} = wP_{\mathrm{hard}}^\top
  \]
  \[
    y_q = Q_h(x_{\mathrm{perm}})\,Q_h(w_{\mathrm{perm}})^\top
  \]

  但反向梯度通过 group surrogate 传播。当前目标写为：
  \[
    \Loss
    = \underbrace{\MSE(y_q, y_{\mathrm{ref}})}_{\Loss_{\mathrm{recon}}}
    + \lambda_{\mathrm{grp}}
    \underbrace{\left(
      \|A_{\mathrm{soft}}\mathbf{1}-\mathbf{1}\|_2^2
      +
      \|A_{\mathrm{soft}}^\top\mathbf{1}-k\mathbf{1}\|_2^2
    \right)}_{\Loss_{\mathrm{group}}}
  \]

  这对应于：
  \begin{itemize}
    \item 每个 channel 的总分配质量为 1
    \item 每个 group 的总容量接近 $k$
  \end{itemize}
\end{frame}

\begin{frame}{9. 训练与评估口径}
  \textbf{训练口径}
  \[
    \Loss_{\mathrm{train}}
    = \Loss\bigl(A_{\mathrm{ste}}\bigr)
  \]

  \textbf{评估口径}
  \[
    \Loss_{\mathrm{hard}}
    = \MSE\!\left(
      Q_h(xP_{\mathrm{hard}}^\top)\,Q_h(wP_{\mathrm{hard}}^\top)^\top,\,
      y_{\mathrm{ref}}
    \right)
  \]

  \vspace{0.4em}
  baseline 则取 identity grouping：
  \[
    P = I
    \quad \Longrightarrow \quad
    \Loss_{\mathrm{base}} = \Loss_{\mathrm{hard}}(I)
  \]

  相对收益定义为
  \[
    \mathrm{Improve}
    = \frac{\Loss_{\mathrm{base}}-\Loss_{\mathrm{final}}}
    {\max(|\Loss_{\mathrm{base}}|,\epsilon)}
  \]

  \vspace{0.4em}
  训练中真正关心的不是单调下降，而是最终 hard 口径是否优于 baseline。
\end{frame}

\begin{frame}{10. 温度退火与 Early Stop}
  温度采用指数退火：
  \[
    \tau_t
    =
    \tau_{\mathrm{start}}
    \left(
      \frac{\tau_{\mathrm{end}}}{\tau_{\mathrm{start}}}
    \right)^{\frac{t}{T-1}}
  \]

  \vspace{0.4em}
  为降低离散训练的高频波动，对 hard eval loss 做 EMA：
  \[
    m_t =
    \begin{cases}
      \ell_t, & t=0\\
      \beta m_{t-1} + (1-\beta)\ell_t, & t\ge 1
    \end{cases}
  \]
  其中 $\ell_t=\Loss_{\mathrm{hard}}^{(t)}$。

  \vspace{0.4em}
  若连续若干轮没有达到相对改善阈值，则 early stop。

  \begin{block}{作用}
    这一步主要是稳住停止条件，而不是保证训练曲线单调。
  \end{block}
\end{frame}

\begin{frame}{11. 局部离散 Refinement}
  训练结束后，从当前最优分组对应的 hard 排列出发，做少量 pair-swap hill climbing：
  \[
    \pi \rightarrow \pi'
  \]
  其中 $\pi'$ 由交换 $\pi$ 中两个位置得到。

  对每个候选 permutation，重新计算
  \[
    \Loss_{\mathrm{hard}}(\pi')
  \]
  若
  \[
    \Loss_{\mathrm{hard}}(\pi') < \Loss_{\mathrm{hard}}(\pi)
  \]
  则接受更新。

  \vspace{0.4em}
  这一步的意义是：
  \begin{itemize}
    \item 训练阶段优化的是可导近似目标
    \item refinement 直接在真实 hard 分组目标上做局部离散搜索
    \item 常能额外捡到一部分 hard loss 收益
  \end{itemize}
\end{frame}

\begin{frame}{12. 多卡并行策略}
  当前实现采用 expert 级并行，而非梯度同步并行。

  对于总 expert 数 $E$、进程数 $W$、当前 rank $r$：
  \[
    \mathcal{E}_r = \{\, r,\ r+W,\ r+2W,\ \dots \,\}
  \]

  每个进程只处理集合 $\mathcal{E}_r$ 中的 expert。

  \vspace{0.5em}
  设备分配：
  \[
    \text{device} = \texttt{npu:LOCAL\_RANK}
  \]

  这种方式的特点：
  \begin{itemize}
    \item expert 间相互独立，天然适合任务并行
    \item 实现比 DDP 更简单
    \item 无需在 rank 间同步 score 参数
  \end{itemize}
\end{frame}

\begin{frame}{13. 结果产物设计}
  每个 expert 会保存一份结果对象，包含：
  \[
    \{\text{config},\ \text{baseline},\ \text{best\_perm},\ \text{best\_score},\ \text{loss history},\ \dots\}
  \]

  关键字段包括：
  \begin{itemize}
    \item $\Loss_{\mathrm{base}}$
    \item $\Loss_{\mathrm{best\_before\_refine}}$
    \item $\Loss_{\mathrm{refined}}$
    \item $\Loss_{\mathrm{final}}$
    \item best group / final best group / best score
    \item loss / EMA / tau / lr 的历史序列
  \end{itemize}

  同时提供后处理脚本，为每个 expert 绘制：
  \begin{itemize}
    \item hard eval loss 曲线
    \item baseline 水平线
    \item refine 前后最优值对比
  \end{itemize}
\end{frame}

\begin{frame}{14. 当前现象与理解}
  当前实验已经实现：
  \begin{itemize}
    \item 方法链路打通
    \item 多卡运行
    \item 结果保存与自动画图
    \item 训练语义从 full permutation 切到 group assignment
  \end{itemize}

  但目前观察到：
  \[
    \Loss_{\mathrm{hard}}^{(t)}
  \]
  在不少 expert 上存在较强波动，整体下降趋势也不总是明显。

  \vspace{0.4em}
  这与以下因素有关：
  \begin{itemize}
    \item grouping 本身仍由排序诱导，离散性依然存在
    \item 训练中同时用了 OT 松弛与 quant STE
    \item soft group surrogate 与 hard 分组评估之间仍存在 gap
    \item 学习率、温度退火等超参可能仍需继续调
  \end{itemize}
\end{frame}

\begin{frame}{15. 下一步工作}
  后续重点并不在工程链路，而在训练稳定性与收益验证：

  \vspace{0.4em}
  \begin{block}{调参与训练稳定性}
    调整学习率、温度终值、退火速度，观察 hard eval 口径是否更稳定优于 baseline。
  \end{block}

  \begin{block}{统计性验证}
    统计更多层、更多 expert 上的收益分布，而不是只看单个曲线。
  \end{block}

  \begin{block}{目标对齐}
    进一步分析 soft 训练路径与 hard 最终目标之间的差距，以及 refine 的贡献占比。
  \end{block}
\end{frame}

\begin{frame}{总结}
  \begin{enumerate}
    \item 当前方案本质上是在 expert 级别学习量化友好的 channel grouping
    \item 方法上保留了 Sinkhorn/OT 的位置松弛，再通过块状矩阵聚合成 group assignment
    \item 工程上已经支持多卡并行、group 结果保存与自动可视化
    \item 当前主要瓶颈仍是训练稳定性与最终收益强度
  \end{enumerate}
\end{frame}

\begin{frame}
  \centering
  {\LARGE 谢谢}
\end{frame}

\end{document}
